% ============================================================================
% CHAPTER 3 - REQUIREMENTS, IMPACTS AND CONSTRAINTS
% Rubric mapping:
%   CO2  (5)  -> 3.1 Final Specifications and Requirements
%   CO3  (5)  -> 3.2 Societal Impact
%   CO4  (5)  -> 3.3 Environmental Impact
%   CO12 (5)  -> 3.4 Ethical Issues
%   CO10 (3)  -> 3.6 Project Management Plan, 3.7 Risk Management
%   CO11 (2)  -> 3.8 Economic Analysis
% This chapter carries 25 of the 100 marks, all from the supervisor.
% ============================================================================

\section{Final Specifications and Requirements}

This section states what the system has to do, what it has to be, and what it
needs in order to run. The requirements were fixed before the final build and were
checked against the delivered system at the end of the project.

\subsection{Functional requirements}

Table \ref{tab:fr} lists what the system must do. Each requirement has a short
identifier so that later chapters can refer to it, and the last column says where
the requirement is met and where it is evaluated.

\begin{table}[H]
\centering
\caption{Functional requirements.}
\label{tab:fr}
\small
\begin{tabular}{p{1.0cm}p{9.2cm}p{3.3cm}}
\toprule
\textbf{ID} & \textbf{Requirement} & \textbf{Where met} \\
\midrule
FR1 & Accept a recorded lecture video in a common container and extract a
16 kHz mono audio track from it. & Chapter 4.5 \\
FR2 & Produce a timestamped transcript of the lecture in romanized Banglish,
using a model adapted to this speech. & Chapter 4.2, result in 5.1 \\
FR3 & Detect the points at which the whiteboard is wiped, and treat the writing
between two wipes as one board. & Chapter 4.3 \\
FR4 & Reconstruct each board without the lecturer in front of it, using only
pixels that appear in some real frame of the video. & Chapter 4.3, result in 5.2 \\
FR5 & Find the separate blocks of writing on a reconstructed board and draw a
numbered coloured box around each one. & Chapter 4.4 \\
FR6 & Name and transcribe the content of each numbered box with a
vision-language model. & Chapter 4.4, result in 5.3 \\
FR7 & Generate a lecture note, one section per board, that refers to the boxes by
number and colour. & Chapter 4.5, result in 5.4 \\
FR8 & Produce the note in two languages, Banglish and English, each written
wholly in its own language. & Chapter 4.5 \\
FR9 & Delete any quotation of the lecturer that does not appear word for word in
the transcript, and any reference to a box number that does not exist. &
Chapter 4.5 \\
FR10 & Keep material that the lecturer did not say inside a separate labelled
section of the note. & Chapter 4.5 \\
FR11 & Output the note as Markdown and as a single self-contained HTML page with
the board images embedded. & Chapter 4.5 \\
FR12 & Record, for every lecture processed, how long each stage took and how many
quotations were kept and deleted. & Chapter 4.6 \\
\bottomrule
\end{tabular}
\end{table}

\subsection{Non-functional requirements}

\begin{table}[H]
\centering
\caption{Non-functional requirements.}
\label{tab:nfr}
\small
\begin{tabular}{p{1.2cm}p{5.0cm}p{7.3cm}}
\toprule
\textbf{ID} & \textbf{Requirement} & \textbf{How it is satisfied} \\
\midrule
NFR1 & Runs on one consumer desktop graphics card & Speech training fits in 12 GB
with gradient checkpointing; the vision-language work fits in 32 GB. No cluster
and no paid cloud service is used. \\
\addlinespace
NFR2 & Reproducible & Every reported number has a command that regenerates it
from files held in the project repository. Splits, seeds and the chosen settings
are written to JSON beside the model. \\
\addlinespace
NFR3 & Restartable & Each stage writes its output to a file and is skipped when
that file already exists, so a run that is interrupted continues instead of
starting again. \\
\addlinespace
NFR4 & No fabricated content in the board image & Every pixel of a reconstructed
board comes from a real frame. Generative inpainting is not used. \\
\addlinespace
NFR5 & Traceable claims in the note & Quotations are checked against the
transcript, board content is passed to the model as text, and non-lecture
material is labelled. \\
\addlinespace
NFR6 & Offline operation & After the model weights are downloaded once, no
component needs the internet. No lecture audio, video or text is sent to any
external service. \\
\addlinespace
NFR7 & Data protection & Videos, transcripts and board images are held in a
private repository. The public output of the project contains no raw recording. \\
\bottomrule
\end{tabular}
\end{table}

\subsection{Data requirements}

The project needed three kinds of data, and the cost of producing them shaped the
whole schedule.

\textbf{Lecture recordings.} Forty-four classroom recordings were collected from
three lecturers, 8.33 hours in total, each a fixed-camera view of a lecturer at a
whiteboard with no slides. The subjects are introductory programming, digital
logic, databases and computer networking.

\textbf{Reference transcripts.} Twenty-eight of the lectures, covering 5.15 hours
of speech in 609 timed segments and about 39,000 words, were transcribed by hand.
Every transcript follows a written transcription standard described in Section
3.5 and validated by a checking script. The remaining sixteen recordings have no
transcript and are used for the vision side and for demonstration only.

\textbf{Board answer keys.} To measure whether board content reaches the notes,
the contents of 45 boards were written out by hand as lists of items, giving 436
items in total. Each item is labelled by kind, which is one of number, name, code,
term or phrase. These keys are used only for scoring and are never given to any
model. A separate human check covered a further 98 boards and recorded, for each
one, whether anything written on the board was missing from the reconstruction.

\subsection{Software and hardware requirements}

\begin{table}[H]
\centering
\caption{Hardware used. The earlier phase of the project ran on a different
machine, and the two are kept separate throughout this report.}
\label{tab:hw}
\small
\begin{tabular}{p{3.0cm}p{2.4cm}p{8.2cm}}
\toprule
\textbf{Machine} & \textbf{GPU} & \textbf{What it was used for} \\
\midrule
Development desktop & RTX 3060, 12 GB & All speech work: data preparation,
hyperparameter tuning, LoRA training, evaluation, board reconstruction and the
person segmentation network. \\
\addlinespace
Second desktop & RTX 5090, 32 GB & All Qwen work: board reading with the
vision-language model, box naming, and note generation for every lecture. \\
\addlinespace
Earlier phase machine & RTX 3090, 24 GB & The first version of the pipeline,
whose results are reported separately where they are still used. \\
\bottomrule
\end{tabular}
\end{table}

The software stack is Python 3.11 on the development desktop with PyTorch
2.5.1 built for CUDA 12.1, and Python 3.12 on the second desktop with PyTorch
2.11 built for CUDA 12.8, which the newer card requires. Both machines run
Transformers 5.12.1 and PEFT for the LoRA adapters. FFmpeg extracts audio and
frames. The models used are whisper-large-v3-turbo for speech,
Qwen2.5-VL-7B-Instruct and Qwen2.5-VL-3B-Instruct for board reading,
Qwen2.5-7B-Instruct for note writing, a pretrained DeepLabV3 with a ResNet-101
backbone for person segmentation \cite{chen2017deeplabv3}, and a WavLM speaker
embedding model used only to confirm which lecturer appears in which recording.
About 31 GB of model weights and roughly 75 GB of video have to be stored.

\subsection{Constraints}

Four constraints were fixed and could not be traded away.

\textbf{Data.} The corpus is what three lecturers agreed to record and what the
team could transcribe by hand in the time available. It is 5.15 hours of
transcribed speech, not tens of hours, and this bounds how strong any speech
result can be.

\textbf{Memory.} Training the 809 million parameter speech model on a 12 GB card
is only possible with gradient checkpointing, which trades computation for
memory. The largest vision-language models were out of reach entirely, since
Qwen2.5-VL-32B is 68 GB of weights and the 72B version is 147 GB.

\textbf{Time.} The whole project runs to a fixed submission date, so a run that
takes fourteen hours can be started once, not five times.

\textbf{Honesty.} A self-imposed constraint, and the one that removed the most
work. No number enters this report unless a command in the project regenerates
it. An earlier version of the project carried several figures whose numbers had
been typed rather than computed, and those were removed rather than defended.

\section{Societal Impact}

\subsection{Who benefits}

The immediate beneficiaries are students in classrooms where teaching is
code-mixed. A student who missed a class, who is slower at English, or who simply
wants to revise, currently has a video and nothing else. The system turns that
video into a document with the board in it. Because the note exists in Banglish as
well as English, a student who thinks in the mixed language is not forced to read
a translation of their own lecture.

There is a second group of beneficiaries who are easy to overlook. Lecturers get a
record of every board they wrote. Nothing else in a whiteboard class preserves
that, and it costs the lecturer no extra work.

\subsection{Language and fairness}

Language technology serves widely spoken standard languages well and serves mixed
speech badly. A student whose classroom language is Bengali-English code-mixed is
poorly served by both English tools and Bengali tools. Treating that mixture as a
first-class target, with its own transcription standard and its own written form,
is a small step against that imbalance. The same argument applies to the choice of
the Roman alphabet, which is what students actually write in.

\subsection{Accessibility}

A text transcript and a text version of the board make a lecture reachable by a
screen reader and searchable by keyword, which a video is not. The system was not
designed as an assistive technology and has not been tested with users who have a
disability, so this is stated as a plausible benefit rather than a demonstrated
one.

\subsection{Risks and harms}

Three risks are worth naming plainly.

\textbf{Wrong notes can mislead.} The generated notes are not guaranteed to be
correct. Chapter 5 gives a real example in which the model wrote that the two
universal gates are XOR gates when the lecturer and the board both said NAND and
NOR. A student who trusts a note like that will learn the wrong thing. Any
deployment must present the notes as machine-generated study aids to be checked
against the lecture, not as authoritative material, and the system keeps the board
image beside the text so that the reader can check.

\textbf{Privacy of the classroom.} A lecture recording captures a person at work
and may capture student voices. This project used recordings made with the
lecturers' permission and kept them private. Any wider use needs consent from the
people recorded and a clear rule about who may see the output.

\textbf{Effects on attendance and on note taking.} Good automatic notes may reduce
attendance, and the act of taking notes by hand has a learning value of its own.
This is a genuine trade-off and is not resolved by this thesis. The honest
position is that the system is aimed at revision and at students who could not
attend, not at replacing attendance.

\subsection{Health, safety, legal and cultural aspects}

There is no physical safety dimension to this work, since nothing is built that a
person interacts with physically. On the legal side, lecture recordings and
whiteboard content belong to the lecturer and the university, so the recordings
are not redistributed and the model weights released by others are used under
their own licences. On the cultural side, the system deliberately preserves the
lecturer's own mixed speech rather than converting it into standard English, which
respects how the class is actually taught.

\section{Environmental Impact}

\subsection{Why this section is short but not empty}

Training machine learning models consumes electricity, and the field has been
criticised for not reporting it \cite{strubell2019energy}. The headline numbers in
that literature come from very large models trained on clusters. The work in this
thesis is at the other end of the scale, and the point of this section is to give
the actual figure rather than to claim the question does not apply.

\subsection{Compute actually used}

All training is parameter-efficient, so no full model was ever trained from
scratch or fine-tuned end to end. Table \ref{tab:compute} lists the machine time
taken by the jobs that produced the results in Chapter 5, read from the job logs.

\begin{table}[H]
\centering
\caption{Machine time for the jobs behind the reported results, from the job
logs. The two rows marked as approximate are added from run records rather than a
single log line.}
\label{tab:compute}
\small
\begin{tabular}{p{7.4cm}p{2.4cm}p{3.6cm}}
\toprule
\textbf{Job} & \textbf{Machine} & \textbf{Time} \\
\midrule
Hyperparameter tuning, rehearsal, 7 trainings and evaluations & RTX 3060 &
about 5.4 h \\
Hyperparameter tuning on the full corpus, 7 trainings and evaluations & RTX 3060 &
6.0 h \\
Final model, seed 42 & RTX 3060 & 1.3 h \\
Final model, seed 1 & RTX 3060 & 1.1 h \\
Second learning rate, two seeds & RTX 3060 & about 2.4 h \\
Transcripts and board reconstruction for all lectures & RTX 3060 &
about 2 h \\
Board reading, box naming and note generation for all lectures & RTX 5090 &
about 6 h \\
\midrule
\textbf{Total} & & \textbf{about 24 h} \\
\bottomrule
\end{tabular}
\end{table}

\subsection{Energy and emissions}

The figures below are an estimate built from the measured hours above and stated
assumptions, not a meter reading.

Assuming a whole-system draw of about 250 W for the development desktop while
training and about 600 W for the second desktop, the 18 hours on the first machine
and 6 hours on the second come to roughly 4.5 and 3.6 kilowatt hours, so about
8 kWh in total. At a grid emission factor of about 0.65 kilograms of carbon
dioxide equivalent per kilowatt hour, which is a reasonable figure for the
Bangladesh grid, that is roughly 5 kilograms of carbon dioxide equivalent for the
entire experimental programme. For comparison, the widely cited estimate for
training a single large language model with neural architecture search is over
280,000 kilograms \cite{strubell2019energy}.

\subsection{Design choices that reduced the footprint}

Three decisions kept the figure small, and they were taken for other reasons as
well.

Using LoRA rather than full fine-tuning means each training run updates a few
million parameters instead of 809 million, which is why a run finishes in about
an hour on a mid-range card. Choosing hyperparameters by a pre-registered plan
with a fixed number of stages meant seven training runs rather than an open-ended
search. Deciding not to run the 32 billion and 72 billion parameter
vision-language models avoided both a 68 GB download and the compute to use them,
and the 3B against 7B comparison reported in Chapter 5 shows that the larger model
was not where the important difference lay.

\subsection{Sustainability of the solution in use}

Once trained, the system runs on one desktop. Processing a lecture end to end
takes a few minutes of graphics card time, so a department could run it on
hardware it already owns. There is no requirement for a data centre, and no
per-lecture payment to an external service, which also means no lecture content
leaves the institution.

\section{Ethical Issues}

\subsection{Consent and the people in the recordings}

Every recording used here was made with the knowledge and permission of the
lecturer in it. Three lecturers took part. The recordings are not published and
are not shared outside the team. Where a board image appears in this report, the
lecturer has been removed by the reconstruction itself, so the figures show
writing rather than people.

\subsection{Data handling}

Videos, audio, transcripts and board images are kept in a private repository. The
transcripts contain classroom speech only. No student personal data was collected.
No lecture content is sent to an external service at any point in the pipeline,
because every model runs locally, which was one of the reasons for preferring open
model weights over a commercial interface.

\subsection{Honesty about results, and how it was enforced}

The most serious professional risk in a project of this kind is reporting a number
that was never computed. This project had that problem in its earlier phase. A
significance value, an effect size, an ablation table and a set of per-video
statistics had been written into figures and text without any code behind them.
They were found during a review, they are listed in the project record as retired
claims, and none of them appears in this report.

Two rules were then adopted and kept. First, a number may be written in the report
only if a command in the project regenerates it, and the command is recorded
beside the number. Second, results that go against the project's own hypothesis
are reported in the same detail as results that support it. Chapter 5 contains
several of the second kind, including a display clean-up that made board reading
slightly worse, an alignment idea that did not validate, and a null result on
whether the improved transcript improves the notes.

\subsection{Risks created by the system itself}

A note generator can state things the lecturer never said. The design contains
three specific controls for this, described in Chapter 4, and Chapter 5 reports
how well each one held, including the one case where a box reference survived the
checker and the counted number of quotations that had to be deleted. The controls
reduce the risk and do not remove it, and the report says so.

There is also an academic integrity dimension. Notes produced by a machine from a
lecture are a study aid. If a student submits them as their own writing, that is
plagiarism in the ordinary sense, and the labelling of generated material inside
the notes makes the origin visible.

\subsection{Professional responsibility}

The project follows the ordinary expectations of engineering practice, which in
this context mean citing the models and datasets that others released, obeying
their licences, not overstating what was built, and keeping the limitations
visible in the same document as the results rather than in a footnote.

% TODO (Rafi): if the department requires a declaration about the use of AI tools
% during writing, put the exact wording here and in core/ethics_statement.tex.

\section{Standards}

The project follows the standards and conventions listed in Table
\ref{tab:standards}. Several of them are conventions of the field rather than
formal standards, and the table says which is which.

\begin{table}[H]
\centering
\caption{Standards and conventions followed.}
\label{tab:standards}
\small
\begin{tabular}{p{3.4cm}p{4.0cm}p{6.2cm}}
\toprule
\textbf{Area} & \textbf{Standard} & \textbf{How it is applied} \\
\midrule
Transcription & Project transcription guide, version 1.0 &
Segments of 10 to 25 seconds and never over 30, timestamps written as
\texttt{[M:SS-M:SS]}, fillers kept, ASCII only, a fixed spelling table for common
words. A validation script checks every file against it. \\
\addlinespace
Audio & 16 kHz mono PCM WAV &
The input format the speech model expects, produced by FFmpeg from any container.
\\
\addlinespace
Speech evaluation & Word and character error rate &
The standard measures of the field, reported as per-clip medians with paired
rank tests, following the practice discussed in Chapter 2. \\
\addlinespace
Text encoding & UTF-8 &
All transcripts, notes and result files. \\
\addlinespace
Interchange formats & JSON, Markdown, HTML5 &
Machine-readable results in JSON, notes in Markdown, and a self-contained HTML
page with images embedded so that a note can be opened with no other files. \\
\addlinespace
Code & PEP 8 for Python, Git for version control &
One script per stage, no deletion of experimental scripts, every result
traceable to a commit. \\
\addlinespace
Referencing & IEEE style &
As required by the department template, produced with biblatex. \\
\addlinespace
Licensing & Licences of the released models &
Whisper, Qwen and the segmentation network are used under the terms their
publishers set, and are cited. \\
\bottomrule
\end{tabular}
\end{table}

\section{Project Management Plan}

\subsection{Phases}

The work ran in five phases.

\textbf{Phase 1, problem and baseline.} The earlier version of the pipeline was
built and measured, which established that the speech side was the binding
problem and that the note quality had never been evaluated against anything.

\textbf{Phase 2, review and correction.} Every claim of the earlier phase was
checked against code. Claims with no computation behind them were retired, and a
project-wide rule was adopted that a number needs a command.

\textbf{Phase 3, data.} The transcription standard was written, the recordings
were collected, and the transcripts were produced and checked word by word. This
phase was the longest in human hours and ran in parallel with the others.

\textbf{Phase 4, experiments.} Speech fine-tuning with its pre-registered tuning
plan, board reconstruction, board reading, and note generation, with evaluation at
each stage. Work was split across two machines, with all speech work on one and
all vision-language work on the other, so that neither waited for the other.

\textbf{Phase 5, writing.} Figures, tables and this report.

\figplaceholder{fig-3-1-gantt}{Project schedule.}{A Gantt chart with one row per
phase, drawn from the list in Section 3.6.1. Rows: Problem and baseline, Review
and correction, Data collection and transcription, Experiments (with the speech
and vision tracks as two sub-rows running in parallel), Writing. Mark the two
fixed deadlines as vertical lines: draft submission and slide submission. Keep
the month names on the horizontal axis.}{7cm}

\subsection{Division of work and resources}

Three students worked on the project under a supervisor and a co-supervisor. The
resources were the two desktops described in Section 3.1.4, both already owned,
the open model weights, and the team's own time. No budget was spent on cloud
computing or on paid annotation.

Two practices mattered more than any tool. Long jobs were written as scripts that
run unattended, commit their own results and push them, so that a fourteen-hour
training run cost no attention. All planning lived in two files in the repository,
one holding the ordered task list and one holding every result with the command
that produces it, so that work could move between machines without anything being
carried in a person's memory.

\subsection{Schedule control}

The fixed points were the draft submission and the slide submission. Work was
scheduled backwards from them. When the corpus was closed at 5.15 hours, the
planned six-hour and ten-hour training runs were cancelled rather than delayed,
because a result that arrives after the deadline is worth nothing. The decision
and the reason were recorded at the time.

\section{Risk Management}

Table \ref{tab:risks} lists the risks identified for the project, what was done
about each, and what actually happened. Several of these risks occurred, so the
last column is a record rather than a prediction.

\begin{table}[H]
\centering
\caption{Risk register, with outcomes.}
\label{tab:risks}
\small
\begin{tabular}{p{3.0cm}p{1.3cm}p{4.4cm}p{4.9cm}}
\toprule
\textbf{Risk} & \textbf{Severity} & \textbf{Mitigation planned} &
\textbf{What happened} \\
\midrule
Not enough transcribed speech &
High &
Use parameter-efficient fine-tuning, which needs less data; report both an
unseen-lecturer design and an unseen-lecture design &
Occurred. The corpus closed at 5.15 h instead of the hoped ten. Both designs are
reported, and the limitation is stated. \\
\addlinespace
Speaker leakage between training and test &
High &
Sort recordings by lecturer and confirm with speaker verification &
Occurred. One recording had been labelled as the wrong lecturer, so an early
result was optimistic. It was found by voice verification, the result was
recomputed and the old number was retired. \\
\addlinespace
Training instability &
High &
Run two random seeds for every reported configuration &
Occurred. At the tuned learning rate one seed collapsed into repetition loops.
Both seeds are reported, and a stable configuration is reported beside them. \\
\addlinespace
Graphics memory exhausted &
Medium &
Gradient checkpointing, smaller batch, second machine for the large models &
Occurred and was handled. Training needed checkpointing to fit in 12 GB. \\
\addlinespace
Model weights too large for the disk &
Medium &
Check free space before any download &
Occurred. The 32 billion parameter vision-language model did not fit and was
dropped, with the reason recorded. \\
\addlinespace
Results that cannot be reproduced &
High &
A command recorded for every number; splits and seeds written to disk &
Controlled. Every table in Chapter 5 has a command behind it. \\
\addlinespace
Generated notes containing false statements &
High &
Ground the notes in board text, fence non-lecture material, check quotations
against the transcript &
Partly controlled. The controls work and are measured, but a factual error in the
generated prose was still found by reading, and it is reported. \\
\addlinespace
Loss of work or of a machine &
Medium &
Version control, results committed and pushed automatically by the jobs
themselves &
Did not occur. \\
\addlinespace
Deadline slip &
High &
Two machines in parallel, unattended overnight jobs, cancel rather than delay &
Controlled. The planned larger runs were cancelled when the corpus closed. \\
\bottomrule
\end{tabular}
\end{table}

\section{Economic Analysis}

\subsection{What the project cost}

The project spent no money on computing. Both desktops were already available, and
every model used is published under an open licence. The real cost is human time,
and most of it is transcription.

\begin{table}[H]
\centering
\caption{Cost of the project. The transcription figure is measured, at about one
hour of human work per ten minutes of video. The electricity figure uses the
assumptions in Section 3.3.3.}
\label{tab:cost}
\small
\begin{tabular}{p{6.0cm}p{3.6cm}p{4.2cm}}
\toprule
\textbf{Item} & \textbf{Quantity} & \textbf{Cost} \\
\midrule
Transcribing 5.15 h of lecture speech by hand & about 31 person-hours &
team time, no cash cost \\
Writing and verifying 45 board answer keys, 436 items & one working day &
team time \\
Checking 98 reconstructed boards for missing writing & 30 to 40 minutes &
team time \\
Graphics card time for every reported experiment & about 24 h & about 8 kWh,
a negligible cash amount \\
Model weights and video storage & about 106 GB & existing disks \\
Graphics hardware & two desktops & already owned \\
\bottomrule
\end{tabular}
\end{table}

\subsection{What it would have cost otherwise}

Two alternatives are worth pricing for comparison.

Renting the same compute is cheap. Twenty-four hours on a rented mid-range
graphics card at typical rates of one to two United States dollars an hour is
about thirty to fifty dollars. Compute is not what makes this kind of project
expensive.

Commercial transcription is the opposite. Bengali speech-to-text services exist,
but they output the Bengali script, so for this purpose they would not remove the
human work at all. Every one of the 31 hours would still be needed, and a fee
would be added on top.

\subsection{Cost and benefit of the system}

The benefit to be weighed is the reduction in human effort to produce a usable
transcript.

Producing a transcript by hand from nothing took about six hours of work per hour
of video in this project, because off-the-shelf output was of no use as a starting
point. The fine-tuned model reaches a median character error rate of 15.8 per
cent, which is a draft that a person can correct rather than retype. If correcting
such a draft takes a third of the time, the cost falls from about six hours per
hour of video to about two. Over the 8.33 hours of recordings collected here, that
is the difference between about fifty hours of work and about seventeen.

This projection has not been tested, and the thesis does not claim it as a result.
Measuring it properly means timing people as they correct real drafts, which is
listed with the reader study in Chapter 6. What can be claimed is the input to the
calculation, which is the measured error rate, and the measured cost of the manual
route.

\subsection{Cost of deployment}

A department that wanted to run this on its own lectures would need one desktop
with a graphics card of 12 GB or more, which is ordinary teaching laboratory
hardware, and no recurring payment. Processing one lecture takes a few minutes of
card time. The recurring cost is storage for the recordings and someone to check
the notes before they are given to students, and that check is the part the system
does not remove.
